% Part: sets-functions-relations
% Chapter: functions
% Section: kinds

\documentclass[../../../include/open-logic-section]{subfiles}

\begin{document}

\olfileid{sfr}{fun}{kin}
\olsection{વિધેયોના પ્રકારો}

\begin{explain}
આપણને વારંવાર મળતા વિધેયોના કેટલાક પ્રકારોનું
વર્ગીકરણ રજૂ કરવું ઉપયોગી થશે.

શરૂઆતમાં એવાં વિધેયો લઈએ જેમના સહપ્રદેશનો દરેક સભ્ય
વિધેયની કિંમત હોય. આવાં વિધેયોને !!{surjective} કહે છે;
તેમને \olref{fig:surjective} પ્રમાણે ચિત્રમાં દર્શાવી શકાય.

\begin{figure}
  \olasset{assets/diagrams/surjective.tikz}
  \caption{!!^a{surjective} વિધેયના સહપ્રદેશનો દરેક
    !!{element} તેની કિંમત હોય છે.}
  \ollabel{fig:surjective}
\end{figure}
\end{explain}

\begin{defn}[!!^{surjective} વિધેય]
વિધેય $f \colon A \rightarrow B$ \emph{!!{surjective}}
હોય જો અને માત્ર જો $B$ એ $f$નો વિસ્તાર પણ હોય, એટલે કે
દરેક $y \in B$ માટે ઓછામાં ઓછો એક $x \in A$
એવો હોય કે $f(x) = y$; અથવા પ્રતીકોમાં:
\[
  (\forall y \in B)(\exists x \in A)f(x) = y.
\]
આવા વિધેયને $A$થી $B$માંનું !!a{surjection} કહીએ છીએ.
\end{defn}

\begin{explain}
$f$ એ !!a{surjection} છે તે બતાવવા, તમારે બતાવવું પડે
કે $f$ના સહપ્રદેશની દરેક વસ્તુ કોઈક આગત $x$ માટે
$f(x)$ની કિંમત છે.

નોંધો કે કોઈ પણ વિધેય !!a{surjection}
\emph{ઉત્પન્ન કરે છે}. કારણ કે આપેલા વિધેય
$f \colon A \to B$ માટે $f'(x) = f(x)$ દ્વારા
$f' \colon A \to \ran{f}$ વ્યાખ્યાયિત કરીએ.
$\ran{f}$ની \emph{વ્યાખ્યા} જ
$\Setabs{f(x) \in B}{x \in A}$ હોવાથી આ વિધેય
$f'$ ચોક્કસ !!a{surjection} છે.
\end{explain}

\begin{explain}
કોઈ પણ વિધેય દરેક સંભવિત આગતને અનન્ય નિર્ગત પર મોકલે છે.
પરંતુ એવાં વિધેયો પણ છે જે જુદાં આગતોને કદી એક જ
નિર્ગત પર મોકલતાં નથી. આવાં વિધેયોને !!{injective}
કહે છે; તેમને \olref{fig:injective} પ્રમાણે ચિત્રમાં
દર્શાવી શકાય.
\begin{figure}
  \olasset{assets/diagrams/injective.tikz}
  \caption{!!^a{injective} વિધેય કદી બે જુદી દલીલોને
    એક જ કિંમત પર મોકલતું નથી.}
  \ollabel{fig:injective}
\end{figure}
\end{explain}

\begin{defn}[!!^{injective} વિધેય]
વિધેય $f \colon A \rightarrow B$ \emph{!!{injective}}
હોય જો અને માત્ર જો દરેક $y \in B$ માટે વધુમાં વધુ
એક $x \in A$ એવો હોય કે $f(x) = y$.
આવા વિધેયને $A$થી $B$માંનું !!a{injection} કહીએ છીએ.
\end{defn}

\begin{explain}
$f$ એ !!a{injection} છે તે બતાવવા, તમારે બતાવવું પડે
કે $f$ના પ્રદેશના કોઈ પણ !!{element}s $x$ અને $y$
માટે, જો $f(x)=f(y)$ હોય, તો $x=y$ થાય છે.
\end{explain}

\begin{ex}
$f(x) = 1$ દ્વારા આપેલું અચળ વિધેય
$f\colon \Nat \to \Nat$ !!{injective} પણ નથી
અને !!{surjective} પણ નથી.

$f(x) = x$ દ્વારા આપેલું તદેવ વિધેય
$f\colon \Nat \to \Nat$ !!{injective} અને
!!{surjective} બંને છે.

$f(x) = x+1$ દ્વારા આપેલું અનુગામી વિધેય
$f \colon \Nat \to \Nat$ !!{injective} છે,
પરંતુ !!{surjective} નથી.

નીચે મુજબ વ્યાખ્યાયિત વિધેય $f \colon \Nat \to \Nat$:
\[
  f(x) =
  \begin{cases}
    \frac{x}{2} & \text{જો $x$ યુગ્મ હોય} \\
    \frac{x+1}{2} & \text{જો $x$ અયુગ્મ હોય.}
  \end{cases}
\]
!!{surjective} છે, પરંતુ !!{injective} નથી.
\end{ex}

\begin{explain}
ઘણી વાર આપણે !!{injective} અને !!{surjective} બંને
હોય એવાં વિધેયો લેવા માગીએ છીએ. આવાં વિધેયોને
!!{bijective} કહીએ છીએ. તેઓ
\olref{fig:bijective}માં દર્શાવેલા વિધેય જેવાં દેખાય છે.
!!^{bijection}sને ક્યારેક \emph{પરસ્પર એક-એક સંગતતાઓ}
પણ કહે છે, કારણ કે તેઓ સહપ્રદેશના ઘટકોને પ્રદેશના
ઘટકો સાથે અનન્ય રીતે જોડે છે.
\begin{figure}
  \olasset{assets/diagrams/bijective.tikz}
  \caption{!!^a{bijective} વિધેય સહપ્રદેશના ઘટકોને
    પ્રદેશના ઘટકો સાથે અનન્ય રીતે જોડે છે.}
  \ollabel{fig:bijective}
\end{figure}
\end{explain}

\begin{defn}[!!^{bijection}]
વિધેય $f \colon A \to B$ \emph{!!{bijective}} હોય
જો અને માત્ર જો તે !!{surjective} અને !!{injective}
બંને હોય. આવા વિધેયને $A$થી $B$માંનું
(અથવા $A$ અને~$B$ વચ્ચેનું) !!a{bijection} કહીએ છીએ.
\end{defn}

\end{document}
